% !TEX root = ../main.tex

\chapter{总结与展望}

\section{全文工作总结}

本文完成了 MarketMind 的设计、实现和初步实验。这个项目最初关注的是一个比较实际的问题：只看 K 线很难知道市场正在讨论什么，也很难把新闻、研报和网页检索结果快速整理成可用于判断的材料。因此，本文没有把股票预测只做成一个离线分类任务，而是尝试搭建一条从信息获取、摘要整理到模型预测的完整链路。

系统实现中，MarketMind 支持按股票代码抓取近期 comments、news、reports、7 日 K 线和 SerpAPI 补充检索结果。原始文本不会直接全部塞给模型，而是先分来源做 Summary，再把更稳定的摘要和行情信息交给综合预测 Prompt。在线页面也围绕这个流程实现了折叠展示、分组摘要、网页正文抓取和最终趋势预测。这个过程让我比较明显地感受到，金融大模型系统的难点不只在模型调用本身，还在于数据清洗、缓存、去重、上下文组织和输出解析这些工程细节。

训练与实验部分以 Qwen2.5-7B-Instruct 为基础模型。本文先构建 7 日 K 线监督数据，再使用 qwen3-max 生成带分析过程的 CoT 微调数据，并采用 LoRA 进行监督微调。为了评估外部文本信息的作用，本文另外构建了包含 K 线、news 摘要和 report 摘要的多源评测数据，标签为未来 1、3、7 个交易日方向。由于 SerpAPI 的历史结果难以精确回放，在线系统中接入了该模块，但离线评测暂时没有把它纳入统一实验集。

实验结果和预期大体一致：只依赖历史 K 线时，模型的表现接近随机二分类；加入新闻和研报摘要后，未来 1 日和 3 日方向预测分别提升 6.89 和 11.08 个百分点，7 日预测没有提升；经过 LoRA 微调后，7 日 K 线单周期任务准确率提升 2.71 个百分点。微调模型在接入外部信息时最高达到 56.13\% 的预测成功率。这个数字并不高，也不能直接说明模型可用于交易，但它显示出新闻和研报更可能帮助较短周期的方向判断。

本文还实现了 GRPO 多周期方向训练入口。该部分目前更接近可运行的后续优化方案，还没有形成完整的 SFT+RL 准确率结果表。相比单纯写成展望，本文已经把多周期数据、Prompt、reward 解析和多卡训练配置整理到代码和附录中，为后续继续训练留下了具体入口。

\section{主要研究结论}

从本文工作中可以得到三个比较直接的结论。首先，股票短期方向预测不能只依赖历史行情。新闻事件、市场讨论和研报观点虽然噪声很大，但在 1 日和 3 日任务上确实带来了更高的方向命中率。其次，大语言模型更适合承担“信息整理器”的角色。它不一定能稳定预测价格，但可以把多源文本压缩为利好/利多、利空风险和不确定性等较清晰的分析材料。最后，微调对这类任务的价值主要体现在格式和任务边界上。没有微调时，模型更容易出现回答不稳定或格式不一致的问题。

另外一个比较重要的体会是，在线系统和离线评测的目标并不完全相同。离线实验需要固定标签、固定样本和可重复指标，因此仍采用未来 1、3、7 日涨跌方向作为评价口径；在线系统更接近辅助分析工具，最终输出未来趋势预测和解释依据，而不是强制给出三个机械标签。本文后期也对代码和论文做了这一区分。

\section{本文不足}

本文仍有不少限制。第一，实验规模和时间窗口还不够大。股票市场有明显阶段性，本文当前结果还不能证明模型在不同市场风格、不同板块和极端行情下都有效。第二，外部信息实验还不是严格单因素消融。部分实验同时包含模型微调和输入信息变化，后续还需要在同一模型、同一解码参数和同一测试集下分别比较 K 线、新闻、研报和完整多源输入的贡献。第三，文本质量控制比较初步。评论中存在情绪宣泄和重复表达，新闻与研报也会受到采集覆盖和发布时间影响，本文主要依靠时间窗口、截断和摘要处理，还没有加入细粒度来源可信度建模。

输出稳定性也是一个实际问题。多源输入变长后，模型偶尔会偏离格式，导致自动解析失败。仅靠自然语言 Prompt 约束并不稳，后续更适合引入 JSON schema、函数调用式输出或单独分类头。最后，本文验证的是方向预测和信息辅助能力，还没有进入完整交易策略层面。真实交易还涉及仓位、成本、流动性、止损和回撤控制，因此本文结果只能作为智能投研辅助信号，而不能直接等同于可实盘执行的策略。

\section{后续研究展望}

后续工作可以从三个方向继续推进。第一是扩展数据和评测范围，包括更长时间跨度、更多股票池，以及牛市、熊市、震荡市、行业轮动和事件驱动行情下的分层评估。第二是把 SerpAPI 这类实时检索结果沉淀成可回放数据，例如保存历史快照、记录来源和时间，并加入可信度排序与证据引用。这样才能更准确地评估 Web 检索增强本身的价值。第三是继续完善模型输出和训练目标。在线系统可以改用更严格的结构化输出；训练端可以从 SFT 检查点继续运行 GRPO，把预测正确性、理由一致性和风险提示质量共同纳入奖励设计。

总体而言，MarketMind 目前还是一个原型系统，但它已经打通了从多源金融信息获取、摘要整理、模型预测到离线评测的主要流程。本文的实验结果不应被理解为“模型已经能稳定预测股票”，更合理的理解是：在股票短期方向判断中，外部文本信息和领域微调确实能提供一定帮助，而要把这类方法推进到更可靠的智能投研系统，还需要更严格的数据回放、更稳的结构化输出和更长周期的实证检验。
